Se $!A \ident (!B \lif !C)$, então $\pi$ termina em
\[
\AxiomC{}
\RightLabel{$\pi_1'$}
\Deduce$!B, \Gamma \fCenter \Delta, !C$
\RightLabel{\RightR{\lif}}
\UnaryInf$\Gamma \fCenter \Delta, !B \lif !C$
\LeftSubproofLabel{$\pi_1$}
\AxiomC{}
\RightLabel{$\pi_2'$}
\Deduce$!B \lif !C, \Pi \fCenter \Lambda, !B$
\AxiomC{}
\RightLabel{$\pi_2''$}
\Deduce$!C, \Pi \fCenter \Lambda$
\RightLabel{\LeftR{\lif}}
\BinaryInf$!B \lif !C, \Pi \fCenter \Lambda$
\RightSubproofLabel{$\pi_2$}
\RightLabel{\Cut}
\BinaryInf$\Gamma, \Pi \fCenter \Delta, \Lambda$
\DisplayProof
\]
A !!{proof} que termina em \Cut,
\[
\AxiomC{}
\RightLabel{$\pi_1'$}
\Deduce$!B, \Gamma \fCenter \Delta, !C$
\RightLabel{\RightR{\lif}}
\UnaryInf$\Gamma \fCenter \Delta, !B \lif !C$
\LeftSubproofLabel{$\pi_1$}
\AxiomC{}
\RightLabel{$\pi_2'$}
\Deduce$!B \lif !C, \Pi \fCenter \Lambda, !B$
\RightLabel{\Cut}
\BinaryInf$\Gamma, \Pi \fCenter \Delta, \Lambda, !B$
\DisplayProof
\]
tem !!{height} de corte menor que a de~$\pi$. A hipótese indutiva fornece
!!a{proof}~$\pi_1''$ de $\Gamma, \Pi \fCenter \Delta, \Lambda, !B$.
A !!{proof} que termina em \Cut,
\[
\AxiomC{}
\RightLabel{$\pi_1'$}
\Deduce$B!, \Gamma \fCenter \Delta, !C$
\AxiomC{}
\RightLabel{$\pi_2''$}
\Deduce$!C, \Pi \fCenter \Lambda$
\RightLabel{\Cut}
\BinaryInf$!B, \Gamma, \Pi \fCenter \Delta, \Lambda$
\DisplayProof
\]
Tanto seu posto de corte quanto sua !!{height} de corte são menores que o
posto de corte e a !!{height} de~$\pi$; assim, a hipótese de indução fornece
uma \Log{G3c}-!!{proof}~$\pi_2'''$ de $!B, \Gamma, \Pi, \Pi \fCenter \Delta, \Lambda,
\Lambda$. Agora aplique a hipótese indutiva à !!{proof} que termina
em \Cut,
\[
\AxiomC{}
\RightLabel{$\pi_1''$}
\Deduce$\Gamma, \Pi \fCenter \Delta, \Lambda, !B$
\AxiomC{}
\RightLabel{$\pi_2'''$}
\Deduce$!B, \Gamma, \Pi, \Pi \fCenter \Delta, \Lambda,
\Lambda$
\RightLabel{\Cut}
\BinaryInf$!B, \Gamma, \Pi \fCenter \Delta, \Lambda$
\DisplayProof
\]
(ela tem posto de corte $< \cutr{\pi}$) para obter !!a{proof} de $\Gamma,
\Gamma, \Pi, \Pi, \Pi \fCenter \Delta, \Lambda, \Lambda, \Lambda$. Obtemos
a !!{proof} $\pi'$ de $\Gamma, \Pi \fCenter \Delta, \Lambda$
aplicando \olref[seq][inv]{prop:G3c-cont-adm}.
